\Article{362}
His principiis positis facili negotio divisores primos functionis $x^\nu-1$ determinare possumus. Supponimus, omnes eos divisores, qui etiam functionem 
\SourceLeaf{241}
aliquam $x^\nu-1$ dividunt, existente $\nu'<\nu$, iam inventos esse, reliquosque investigare proponi. Hi autem omnes in hac expressione comprehendi possunt $(P,\rho^k)$, si $P$ sit unus ex ipsis et pro $k$ omnes numeri minores quam $\nu$ ad ipsumque primi substituantur.

In Cap. \textsc{vi} ostendimus, quomodo radices primae aequationis $x^\nu=1$ ita in classes discerpi possint, ut, omnibus per alicuius potestates expressis, eadem in classes distributio habeatur, quaecunque radix prima pro hac basi accipiatur; \emph{periodos} huiusmodi radicum complexus vocavimus. Iam patet, functiones $x,x^\alpha,x^\beta,x^\gamma$ etc., designantibus $\alpha,\beta,\gamma$ etc. omnes numeros ad $\nu$ primos, simili modo in periodos resolvi posse, quamque periodum maiorem rursus in minores donec tandem ad periodos formae $x^k,x^{kp},x^{kpp}\dotstwo x^{kp^{m-1}}$ perveniatur. Hoc ita facto patet

$1^{\mathrm{o}}$ Quoniam periodus quaeque ex huiusmodi periodis minimis $x^k+x^{kp}+\text{etc.}$ composita est, si per quamcunque functionem primam $m$ dimensionum dividatur, residuum fore numerum.

$2^{\mathrm{o}}$. Quum omnes periodi termini semper ad hanc formam reduci queant $x^{\kappa\pdot a^\alpha b^\beta c^\gamma}\dotstwo$, ubi $\kappa.\ a,b,c\dotstwo$ sunt numeri determinati, pro $\alpha,\beta,\gamma\dotstwo$ autem omnes valores substitui possunt; patet, periodum in se ipsam mutari, si pro $x$ substituatur $x^k$ et $k$ sit formae $a^\alpha b^\beta c^\gamma\dotstwo(\mathrm{mod.}\,\nu)$, unde facile perspicitur omnes functiones $P,(P\;\rho^k)$ etc., designante $k$ huiusmodi numerum, si periodus per eas dividatur, idem residuum dare.

$3^{\mathrm{o}}$. Quare periodus subducto tali residuo per productum ex omnibus functionibus $(P,\rho^k)$ dividi poterit.
